\section{Code Listings}

\begin{lstlisting}[language=bash,caption={Automation script for benchmark runs},
label={lst:run-measurements}]
#!/bin/bash

set -e

BENCH="bench"
RUNS=$1
SLEEP_BETWEEN=5
LOGDIR="./logs"

if [ -z "$RUNS" ]; then
    echo "Usage: $0 <num_runs>"
    exit 1
fi

mkdir -p "$LOGDIR"
GOVFILE="/tmp/orig_governors.txt"

sudo touch $GOVFILE
sudo chown root:root $GOVFILE

> "$GOVFILE"
for cpu in /sys/devices/system/cpu/cpu[0-9]*; do
    gov=$(cat "$cpu/cpufreq/scaling_governor")
    echo "$cpu $gov" >> "$GOVFILE"
done

for cpu in /sys/devices/system/cpu/cpu[0-9]*; do
    echo performance | sudo tee "$cpu/cpufreq/scaling_governor" >/dev/null
done

for i in $(seq 1 $RUNS); do
    ts=$(date +"%Y-%m-%d_%H-%M-%S")
    logfile="$LOGDIR/run_${i}_$ts.log"
    echo "[RUN $i/$RUNS] $(date)" | tee -a "$logfile"
    $BENCH | tee -a "$logfile"
    sleep $SLEEP_BETWEEN
done

while read -r cpu gov; do
    echo $gov | sudo tee "$cpu/cpufreq/scaling_governor" >/dev/null
done < "$GOVFILE"
\end{lstlisting}
\clearpage

\begin{lstlisting}[language=c,caption={\texttt{monitor.c}},
label={lst:monitor}]
#include <cpuid.h>
#include <signal.h>
#include <stdio.h>
#include <stdlib.h>
#include <unistd.h>
#include <time.h>

#include <libsmu.h>
#include <pm_table.h>

volatile static int running = 1;
static smu_obj_t smu;

#include <utils.h>

void handle_signal(int sig)
{
    switch(sig) {
        case SIGINT:
        case SIGABRT:
        case SIGTERM:
            running = 0;
        default:
            break;
    }
    // Re-enable the cursor.
    fprintf(stdout, "\e[?25h");
}

int main()
{
    
    if (getuid() != 0 && geteuid() != 0) {
        fprintf(stderr, "Program must be run as root.\n");
        return 1;
    }

    smu_return_val ret = smu_init(&smu);
    if (ret != SMU_Return_OK) {
        fprintf(
            stderr,
            "Error initializing userspace library: %s\n",
            smu_return_to_str(ret)
        );
        return 1;
    }

    if (!smu_pm_tables_supported(&smu))
    {
        fprintf(stderr, "Error: pm_tabes are not supported.\n");
        return 1;
    }

    unsigned int cores = get_cores_count();

    unsigned char* pm_buf = calloc(smu.pm_table_size, sizeof(unsigned char));
    ppm_table_0x240903 pmt = (ppm_table_0x240903)pm_buf;

    int ticks = 10;

    while(running && ticks)
    {
        ret = smu_read_pm_table(&smu, pm_buf, smu.pm_table_size);
        if (ret != SMU_Return_OK)
        {
            fprintf(
                stdout,
                "failed to read pm_table: %s\n",
                smu_return_to_str(ret)
            );
            continue;
        }

        fprintf(stdout, "\e[1;1H\e[2J");
        fprintf(stdout, "Package Power: %f W\n", pmt->SOCKET_POWER);
        fprintf(stdout, "Core Power: %f W\n", pmt->VDDCR_CPU_POWER);
        if (cores != -1)
        {
            for(int i = 0; i < cores; i++)
            {
                fprintf(
                    stdout, "Core[%d] Power: %f W\n", i, pmt->CORE_POWER[i]
                );
            }
        }
        fprintf(stdout, "\e[?25l");
        fflush(stdout);
        struct timespec t, _;
        t.tv_sec = 1;
        t.tv_nsec  = 0;
        nanosleep(&t, &_);
    }

    smu_free(&smu);
}
\end{lstlisting}
\clearpage

\begin{lstlisting}[language=c,caption={\texttt{profile.c}},
label={lst:profile}]
#include <cpuid.h>
#include <signal.h>
#include <stdio.h>
#include <stdlib.h>
#include <unistd.h>
#include <time.h>

#include <libsmu.h>
#include <pm_table.h>

volatile static int running = 1;
static smu_obj_t smu;

#include <utils.h>

void handle_signal(int sig)
{
    switch(sig) {
        case SIGINT:
        case SIGABRT:
        case SIGTERM:
            running = 0;
        default:
            break;
    }
}

int main()
{
    
    if (getuid() != 0 && geteuid() != 0) {
        fprintf(stderr, "Program must be run as root.\n");
        return 1;
    }

    smu_return_val ret = smu_init(&smu);
    if (ret != SMU_Return_OK) {
        fprintf(
            stderr,
            "Error initializing userspace library: %s\n",
            smu_return_to_str(ret)
        );
        return 1;
    }

    if (!smu_pm_tables_supported(&smu))
    {
        fprintf(stderr, "Error: pm_tabes are not supported.\n");
        return 1;
    }

    unsigned int cores = get_cores_count();

    unsigned char* pm_buf = calloc(smu.pm_table_size, sizeof(unsigned char));
    ppm_table_0x240903 pmt = (ppm_table_0x240903)pm_buf;

    int ticks = 10;

    while(running && ticks)
    {
        ret = smu_read_pm_table(&smu, pm_buf, smu.pm_table_size);
        if (ret != SMU_Return_OK)
        {
            fprintf(
                stdout,
                "failed to read pm_table: %s\n",
                smu_return_to_str(ret)
            );
        }
        fprintf(stdout, "Package Power: %f W\n", pmt->SOCKET_POWER);
        fprintf(stdout, "Core Power: %f W\n", pmt->VDDCR_CPU_POWER);
        if (cores != -1)
        {
            for(int i = 0; i < cores; i++)
            {
                fprintf(
                    stdout, "Core[%d] Power: %f W\n", i, pmt->CORE_POWER[i]
                );
            }
        }
        ticks--;
        struct timespec t, _;
        t.tv_sec = 0;
        t.tv_nsec  = 1e6;
        nanosleep(&t, &_);
    }

    smu_free(&smu);
}
\end{lstlisting}
\clearpage
